%% Consignes : « La presentation doit etre structuree autour de la demarche
%% de projet presentee precedemment. [...] Chaque fois que necessaire,
%% indiquer les productions. [...] Il est aussi pertinent d'indiquer les
%% problemes rencontres et les solutions ou strategies mises en oeuvre. »

This section follows the phases of Section~\ref{sec:planning}: the model, the system built
from it, the experiment designed to test it, what the experiment measured, and the problems
met along the way.

\subsection{The Model}
\label{sec:realisation_modele}

\subsubsection{Real urgency}

Real urgency is not read from a single deadline-proximity signal. It is aggregated at each
node from six indicator blocks selected from the literature of
Section~\ref{sec:etatdelart}: time, capacity, performance, risk, cost and carbon. Each
produces a partial urgency $u_m \in [0,1]$, and the aggregation is a weighted probabilistic
OR:

\begin{equation}
    U_{r,\text{local}} = 1 - \prod_{m} (1 - u_m)^{\omega_m}.
    \label{eq:ur}
\end{equation}

The choice of that form over a weighted average is the central modelling decision. A weighted
average is fully compensatory: a node with a completely failed capacity and healthy cost,
carbon and performance readings would score as moderately urgent. Equation~\ref{eq:ur} is
not: any block approaching its maximum drives the aggregate toward the maximum regardless of
the others. For a signal whose purpose is to catch the node that is about to fail,
non-compensatory behaviour is the required property.

\begin{table}[htbp]
\centering
\caption{The six indicator blocks and what each measures}
\label{tab:blocs}
\small
\begin{tabularx}{\linewidth}{>{\raggedright\arraybackslash}p{2.0cm} >{\raggedright\arraybackslash}p{2.7cm} X}
\toprule
\textbf{Block} & \textbf{Role} & \textbf{What it measures} \\
\midrule
Time & Direct cause & Probability that the realistic lead time exceeds the margin to the
next deadline. The only resource whose exhaustion cannot be reversed \\
Capacity & Direct cause & Saturation of storage, weight and throughput, which adds waiting
time a nominal lead time does not show \\
Performance & Direct cause & Drift in the productive health of the resource making the part,
through a composite availability, performance and quality proxy \\
Risk & Direct cause & Expected unavailability: failure probability times recovery time times
severity, revised weekly from declared events \\
Cost & Arbitration & Financial strain. Paying premium freight to protect a deadline is a
symptom, not a cause, and it bounds which responses stay affordable \\
Carbon & Arbitration & Emissions above target. Restricts the admissible set of responses,
typically in the sea-against-air trade-off \\
\bottomrule
\end{tabularx}
\end{table}

Table~\ref{tab:blocs} gives the six and separates them by role. Four describe a mechanism
that leads to rupture; two constrain which responses remain available once it threatens. Each
candidate indicator was screened against four questions before being retained: does it have a
plausible causal link to rupture, can it be measured or credibly estimated, can an operator
act on it, and does it add information the retained blocks do not already carry.
Figure~\ref{fig:blocs} draws the resulting structure.

\begin{figure}[htbp]
    \centering
    \includegraphics[width=0.80\linewidth]{Images/4.Operations/causal_tree_blocks.png}
    \caption{Causal structure behind the six-block selection}
    \label{fig:blocs}
\end{figure}

Three properties of that equation were specified explicitly, because a bare formula does not
state what its bounds mean. At $u_m = 0$ the block is nominal and its factor is $1$, leaving
the product unchanged. At $u_m = 1$ the factor is $0$ and the aggregate saturates. A block
whose data is unavailable is given weight $\omega_m = 0$, which is deliberately distinct from
$u_m = 0$: the first says nothing is known, the second asserts the block is healthy.

\paragraph{The assumption that the blocks are exchangeable.}
The product is symmetric in its factors, so it treats the six blocks as interchangeable up to
their weights. That does not hold cleanly. A value of $u_{\text{time}} = 0.3$ is a
probability of missing a deadline; $u_{\text{cost}} = 0.3$ is an ordinal position on a
normalising scale wearing a cardinal number. They are not equally informative about rupture.
The consequence was carried rather than hidden: the aggregate is used for ranking and for
comparison against the declared value, never as a calibrated probability in its own right.

\subsubsection{Declared urgency and the adequacy score}

Declared urgency is not collected as a free-text estimate. At each node the operator compares
four criteria pairwise through AHP, which yields a consistency-checked weight per criterion,
then rates each on a severity scale. The four criteria were chosen so that each has a
computed counterpart: operational impact and downstream dependencies mirror the severity the
propagation amplifies, the time window mirrors the time block, and recoverability mirrors the
recovery term of the risk block. Without that correspondence the later comparison would be
between unrelated quantities.

\begin{figure}[htbp]
    \centering
    \includegraphics[width=0.76\linewidth]{Images/4.Operations/questionnaire_screenshot.png}
    \caption{The pairwise comparison questionnaire and its consistency gate}
    \label{fig:questionnaire}
\end{figure}

Figure~\ref{fig:questionnaire} shows the questionnaire as the operator meets it. A weight
vector whose consistency ratio exceeds the accepted threshold is refused, and the operator is
sent back to the most contradictory of their own comparisons before the assessment can enter
the network. Without that gate an incoherent judgment would propagate as though it were a
considered one.

Both signals then propagate across the supplier network, and in opposite directions: declared
urgency descends from the customer toward suppliers, real risk ascends from suppliers toward
the customer. The asymmetry is deliberate. Demand pressure is something a customer transmits
to its suppliers; physical risk is something suppliers transmit to their customer. A model
that propagated both the same way would be describing neither.

Figure~\ref{fig:propagation} shows the two signals travelling in opposite directions across
the same graph.

\begin{figure}[htbp]
    \centering
    \includegraphics[width=0.80\linewidth]{Images/4.Operations/urgency_propagation.png}
    \caption{Bidirectional propagation across the supply network}
    \label{fig:propagation}
\end{figure}

Their mismatch, once both are propagated, splits into two named pathologies: false urgency
where the declaration exceeds the state, hidden risk where the state exceeds the declaration.
It is scored on $[0,100]$ through an asymmetric penalty that prices hidden risk at roughly
twice false urgency, following the clinical principle of Section~\ref{sec:pourquoi}.

\begin{figure}[htbp]
    \centering
    \includegraphics[width=0.80\linewidth]{Images/2.Problem/urgency_gap.png}
    \caption{Declared and real urgency, and the two mismatch pathologies}
    \label{fig:gap}
\end{figure}

Figure~\ref{fig:gap} shows the two signals over the life of a task. The step line is the
declaration, which moves only when a review is submitted; the smooth curve is the computed
state, which moves continuously. Where the step line sits above the curve the enclosed area
is false urgency; where the curve sits above the step line it is hidden risk.

\subsection{The System}
\label{sec:realisation_systeme}

The model was delivered as SupplyScore, a standalone Python application of 73 modules across
10 packages, covered by 1\,779 automated test functions in 106 test files, driving a
multi-page web interface over a per-node database layer.

Four architectural decisions carry the engineering weight. Every operation passes through a
single orchestrating facade, so there is no second write path. Data is split into one
registry database plus one database per node, which enforces that one supplier's history is
not readable through another's. Every business write crosses a single mutation service that
validates ranges, diffs against current state and records the change with its author, making
the audit trail a property of the architecture rather than a convention. And each project
carries its own clock, which is what allows a six-month scenario to be replayed in
compressed time without touching the model.

\begin{figure}[htbp]
    \centering
    \includegraphics[width=0.76\linewidth]{Images/4.Operations/hebdo.png}
    \caption{The weekly review page}
    \label{fig:hebdo}
\end{figure}

Figure~\ref{fig:hebdo} shows the operator's main entry point. Its four sections, the
questionnaire, the week's indicator values, milestone progress and any events, commit as a
single signed transaction so that no partial review can enter a node's history.

\begin{figure}[htbp]
    \centering
    \includegraphics[width=0.74\linewidth]{Images/4.Operations/uml_architecture.png}
    \caption{System architecture: single facade, audited writes, per-node persistence}
    \label{fig:archi}
\end{figure}

Figure~\ref{fig:archi} shows how those pieces relate. The single entry point is what makes
the audit guarantee enforceable rather than aspirational: there is no path into the data that
bypasses it, including from the administration interface.

Two capabilities proved more important than anticipated. The explainability page reconstructs
any score term by term, decomposing the declared value into one contribution per criterion
and the computed value into one per block, so that a coordinator asking why a node scored as
it did receives an exact answer rather than a plausible one
(Figure~\ref{fig:explain}). That is what made the diagnosis in
Section~\ref{sec:realisation_resultats} possible from the application rather than by hand.
The export produces frozen data for analysis outside the tool, which the whole experimental
analysis relies on.

\begin{figure}[htbp]
    \centering
    \includegraphics[width=0.78\linewidth]{Images/4.Operations/explainability_waterfall.png}
    \caption{Per-node explainability: exact decomposition of both signals}
    \label{fig:explain}
\end{figure}

\subsection{The Experiment}
\label{sec:realisation_experience}

Building a system that computes a number establishes nothing about whether the number means
anything. The mission therefore included an experiment designed before the system was
finished.

The campaign replays the 2020 to 2022 semiconductor shortage over nineteen turns across an
eight-node satellite supply chain, from a raw wafer supplier through a foundry, a
distributor, a freight operator and an electronics manufacturer to the satellite prime. Real
urgency is fed from public time series of the actual crisis; declared urgency comes from the
participants through the weekly questionnaire. The two signals therefore have independent
sources, which is the precondition for the gap between them to mean anything.

\begin{figure}[htbp]
    \centering
    \includegraphics[width=0.94\linewidth]{Images/4.Operations/helios_network_coefficients.png}
    \caption{The campaign network and its per-arc coefficients}
    \label{fig:reseau}
\end{figure}

Figure~\ref{fig:reseau} gives the network, with the coefficient attenuating demand as it
descends and the one amplifying risk as it ascends, and each node's tier. The foundry at the
deepest tier is the expected bottleneck. The dashed arc into the secondary foundry is
deliberately inert: it represents the announced alternative supplier that in the real crisis
had no spare capacity, and modelling it as present but carrying nothing is the point.

Four artefacts were frozen before any data existed: the scenario with its twelve calibrated
events, the data pack with a hash recorded per file, the role cards, and a register of six
hypotheses each carrying a numeric threshold. Freezing the thresholds in advance is what
prevents a disappointing result from being rescued by moving the bar afterwards.

The design compares two arms on the identical frozen scenario. In the control arm, eight
isolated declarants each see only their own role card and turn briefing and declare blind. In
the treatment arm, the same declarants are additionally shown the system's own predictive
analysis of their node before declaring, and record whether it influenced them. Because the
scenario is frozen, the forecasts are identical between arms, so any behavioural difference
belongs to the act of displaying them.

\subsection{Results}
\label{sec:realisation_resultats}

\subsubsection{The signal does not predict outcomes}

Scored against the outcome definition the system itself applies, meaning a missed milestone,
an abandoned node or a critical event within four weeks, the forecast layer fails on both
axes that matter.

Discrimination is indistinguishable from chance: the area under the ROC curve ranges from
$0.37$ to $0.56$ across horizons and every confidence interval contains $0.5$. That statement
carries a caveat, since the positive class holds between two and six observations, which is
too few to conclude in either direction.

Calibration fails much more clearly, and does not depend on that small sample. The Brier
skill against a forecaster that simply announces the base rate is between $-3.4$ and $-6.9$,
meaning the model is several times worse than the trivial baseline.
Table~\ref{tab:fiabilite} localises the failure.

\begin{table}[htbp]
\centering
\caption{Reliability of the four-week forecast}
\label{tab:fiabilite}
\small
\begin{tabularx}{\linewidth}{l >{\centering\arraybackslash}X >{\centering\arraybackslash}X >{\centering\arraybackslash}X >{\raggedright\arraybackslash}p{3.6cm}}
\toprule
\textbf{Announced band} & \textbf{$n$} & \textbf{Mean announced} & \textbf{Observed} &
\textbf{Reading} \\
\midrule
$[0.0,\ 0.1)$ & 74 & 0.070 & 0.081 & Well calibrated \\
$[0.1,\ 0.2)$ & 24 & 0.128 & 0.000 & None materialised \\
$[0.2,\ 0.5)$ & 1  & 0.354 & 0.000 & None materialised \\
$[0.5,\ 0.9)$ & 7  & 0.859 & 0.000 & None materialised \\
$[0.9,\ 1.0]$ & 14 & 0.992 & 0.000 & None materialised \\
\bottomrule
\end{tabularx}
\end{table}

The pattern is unambiguous, and Figure~\ref{fig:fiabilite} shows it, with each marker's area
proportional to the number of forecasts in its band and the diagonal marking perfect
calibration. Where the model announces a low probability it is right: 74 forecasts averaging
$0.070$ were followed by an adverse outcome $8.1\%$ of the time. Where it announces a high
probability it is wrong every time: twenty-one forecasts announcing a mean of $0.948$ were
followed by the predicted outcome zero times.

\begin{figure}[htbp]
    \centering
    \includegraphics[width=0.84\linewidth]{Images/armA_reliability.png}
    \caption{Reliability of the four-week forecast}
    \label{fig:fiabilite}
\end{figure}

Separating these two failures matters for what happens next, and it is the distinction Murphy
established for exactly this purpose~\autocite{murphy1973vector}. A forecast that orders
outcomes correctly but announces the wrong numbers is repaired by a thin post-hoc layer that
leaves the model untouched. A forecast that does not order outcomes at all needs the model
changed. On this campaign the evidence for the second is weak, resting on very few positive
observations, while the evidence for the first is strong and does not depend on that sample
at all.

\subsubsection{Why: one block explains all of it}

The explainability layer localised the cause. Across all twenty-one confident forecasts, the
four-week probability and the milestone-miss probability never differ by more than $0.020$:
the confident tail is the milestone channel and nothing else. Set against the registry, ten
of the eighteen campaign milestones completed, eight remained active and only one was
actually missed. The model announced near-certain misses for six nodes that all delivered.

Reading the time block directly at the turn the scenario's critical event lands confirms it
from the other side: the block reads numerically zero at six of the eight nodes in the middle
of a crisis. The cause is structural. The block computes milestone risk from a lead-time
distribution against a deadline, while the event engine writes to capacity, cost and failure
probability and never to lead time. No path connects a shock to the quantity the block reads,
so the block cannot respond to a disruption. It is blind by construction, and because it
drives the entire confident tail, so is the forecast.

This is a defect in the model rather than a property of the scenario. Any deployment would
inherit it, and the repair is a change to the event mapping rather than to the mathematics.
It is the single most actionable finding of the mission.

\subsubsection{What happens when the forecast is shown to the operator}

The treatment arm produced a second result, and one that was not sought.

For the first four turns the forecast is hidden and all eight declarants record no influence.
At the turn it first becomes visible, all eight record one, and the direction is entirely
one-sided: five revised their declared urgency downward and three had their existing
assessment confirmed. Not one revised upward. The first upward revision anywhere appears the
following turn, and the scenario's single critical event lands two turns after the moment
every declarant was either reassured or confirmed.

\begin{figure}[htbp]
    \centering
    \includegraphics[width=0.90\linewidth]{Images/armB_influence.png}
    \caption{Influence of the displayed forecast, by turn}
    \label{fig:armB}
\end{figure}

Figure~\ref{fig:armB} shows the distribution turn by turn. Over the full campaign the balance
is close to even, 15 upward against 16 downward, so the one-sided effect belongs to first
contact rather than being a permanent bias.

Two things make this worse than a curiosity. The forecast being displayed comes from the same
layer whose confident predictions never materialise, and whose time block is blind to the
very shock about to arrive. And it is displayed as a numeric probability with a confidence
interval, a format the automation-bias literature identifies as increasing
reliance~\autocite{goddard2012automation,bansal2021whole}. The system was presenting an
unfounded reassurance in its most persuasive available form, immediately before the event it
had failed to anticipate.

One qualification cuts against the strongest reading: the recorded field captures the effect
on the declared level, not on the action taken, and the free-text notes show declarants still
acting protectively. The effect is on the measurement, which is nonetheless the input the
whole adequacy layer consumes.

\subsubsection{The outcome definition moves the answer}

A final measurement concerns how these results should be read at all. Scoring the identical
forecasts against a plausible but naive outcome definition, counting any recorded event
rather than the production definition, moves the base rate from $5.0\%$ to $38.7\%$ and
improves the apparent Brier skill from $-3.36$ to $-0.64$, while leaving discrimination
essentially unchanged.

The model did not change. Only the definition of truth did. An analyst choosing the easier
target would report a system five times better calibrated without a line of code being
touched. A calibration figure quoted without its outcome definition therefore carries little
information, which is a methodological result worth more to the programme than the score
itself.

\subsection{Problems Encountered and How They Were Resolved}
\label{sec:problemes}

Four difficulties shaped the mission, beyond the risks of Section~\ref{sec:risques}.

\paragraph{The first model could not represent the problem.}
The initial formulation expressed urgency as a single time-driven curve. It could not
represent a node on schedule but out of capacity, which is a common and dangerous state. The
resolution was to abandon it in April in favour of multi-block aggregation. The cost was
roughly three weeks; the alternative was building a validated implementation of an
inadequate model.

\paragraph{Measuring the gap requires two independent signals, which is harder than it looks.}
The same difficulty recurred three times in different disguises. A first synthetic dry run
derived the declared value from the computed one, so no gap could appear. The blind arm used
declarants who read their briefings so carefully that their declaration reproduced the
computed state. The treatment arm closed the loop deliberately by showing the model's output
to the declarant. Each time the symptom looked like a modelling problem and each time the
cause was circularity in the measurement. The resolution was to state the independence
requirement as a first-class design constraint and to record, for every result, which of the
two signals could have contaminated the other.

\paragraph{The literature was not accessible from inside the company.}
As set out in Section~\ref{sec:risques}, ALTEN's Direction de l'Innovation holds no journal or
database subscriptions. The state of the art was built entirely on ESTIA and Cranfield
library access. The practical resolution was to front-load literature retrieval into the first
weeks while planning which sources would be needed later, since retrieving a paper was never
a same-hour operation.

\paragraph{Compute was limited to a CPU.}
No GPU was available. This bounded what could be attempted and is the reason one comparison
was run at reduced scale. The resolution was methodological rather than material: choosing
approaches whose cost is dominated by the quantity of data rather than by model size, and
accepting the reduced scale explicitly in the reporting instead of quietly omitting the
comparison.
